%% kapa: exam=1 round=25 year=2014 period=1 form=A subject=economics subject_ko=경제원론 seq=33 number=73
%% source: Q-Net 감정평가사 기출문제 (cst00309) artlSeq=5101207 file=extracted/1008/5101207/2014년 제25회 감정평가사 1차시험 기출문제/1교시(A형).hwp
%% conversion: kordoc 4.24.51 -> tools/gen1.py (auto); review=proofread-sonnet; flags=none
\begin{question}{73}
거시경제의 단기균형과 장기균형에 관한 설명으로 옳은 것은?
\begin{choices}
  \item 물가가 하방경직적일 때 총수요는 단기적으로 실질GDP에 영향을 미친다.
  \item 통화정책과 재정정책은 장기적으로만 실질GDP에 영향을 미친다.
  \item 고전적 이분성은 단기에만 성립하고 장기에는 성립하지 않는다.
  \item 물가와 임금은 장기에 있어서만 경직적이다.
  \item 통화정책은 단기적으로는 명목GDP에 영향을 미치며 장기적으로는 실질GDP에 영향을 미친다.
\end{choices}
\end{question}
\answer{1}
